% Results. Drops into the main paper with % Results. Drops into the main paper with % Results. Drops into the main paper with % Results. Drops into the main paper with \input{results}; also compiles
% standalone (the guard falls back to article when the workshop style is absent).
\documentclass{article}

\IfFileExists{neurips_2026.sty}{\usepackage[dblblindworkshop]{neurips_2026}}{%
  % standalone fallback: lmodern gives microtype a scalable font to expand
  \usepackage[margin=1in]{geometry}\usepackage{lmodern}}

\usepackage[utf8]{inputenc}
\usepackage[T1]{fontenc}
\usepackage{hyperref}
\usepackage{url}
\usepackage{booktabs}
\usepackage{amsfonts}
\usepackage{amsmath}
\usepackage{amssymb}
\usepackage{microtype}
\usepackage{graphicx}
\usepackage{xcolor}

\title{Branchpoint: Towards Structured Multiverse Analysis}

\author{%
  Anonymous Author(s) \\
  Affiliation \\
  \texttt{email} \\
}

\begin{document}
\maketitle

\section{Results}
\label{sec:results}

Eighteen conditions, $n = 50$ analyses each: three models $\times$ two datasets
$\times$ \{Naive, Branchpoint, BpDynamic\}. 896 analyses in total, at a compute
cost of \$145.

\subsection{Branching produces distinct analyses; the allocator decides how many}
\label{sec:results-coverage}

\input{table_coverage_spread}

Table~\ref{tab:coverage-spread} reports the 95\% width of each condition's
Fisher-$z$ distribution in absolute units, against the human panel for the same
dataset, together with what each distinct answer cost.

A second quantity, not tabulated, is needed to read those widths.
\emph{Coverage} is the share of in-scope analyses that are a distinct
standardized endpoint --- how much of what an arm produced is a different answer
rather than the same one again. It rises across the arms for every model and
both datasets: on \texttt{gpt-5.6-sol} it goes $39 \to 68 \to 95$ on blue tit
and $20 \to 69 \to 80$ on Eucalyptus. The naive controls show why it matters.
\texttt{eu\_opus\_naive} produced 50 analyses containing 10 distinct answers,
and \texttt{bt\_opus\_naive} 10 in-scope analyses containing 3. Width alone
reports those as narrow; coverage reports them as \emph{repetitive}, which is a
different failure and the one branching is meant to address.

The allocator does what it is designed to do, and the effect is exact rather
than statistical. The BpDynamic conditions spent 50 branches on 50 distinct
decisions. The Branchpoint conditions, drawing at random from the same open set,
spent 50 branches on 34--37 distinct decisions, wasting the remainder on
decisions already branched.

\paragraph{Covering more decisions does not widen the distribution.}
This is the result we would least like to be true. Across the six cells where
both arms ran, BpDynamic covers more of the decision space in four and is wider
in four --- but they are not the same four. Only \texttt{eu}/\texttt{sol} and
\texttt{eu}/\texttt{opus} win on both. On blue tit BpDynamic leads on coverage
in two of three cells and is the narrower arm in two of three. The additional
decisions the allocator reaches are, on that dataset, ones whose alternatives do
not move the estimand.

\subsection{Nothing converged}
\label{sec:results-convergence}

Every branching condition stopped at \texttt{cap-reached}. Across the twelve of
them, 2{,}157 of 2{,}159 detected decisions remain open: \emph{two} were ever
explored to exhaustion. With $n = 50$ against roughly 200 detected decisions per
condition, even perfect allocation reaches a quarter of them at one alternative
apiece.

Every width reported here is therefore a lower bound from a partial sample, and
the convergence curves are still climbing at $n = 50$. This is why coverage is
worth stating alongside the widths: a width without a denominator invites
reading a 25\%-coverage sample as an estimate of the multiverse.

\subsection{The analyses are far more homogeneous than the humans'}
\label{sec:results-homogeneity}

Under the Fisher-$z$ convention the standard error depends only on the degrees of
freedom, so its variation across a condition measures how much the analyses
differ in effective sample size. The human panels have a standard-error
coefficient of variation of $0.584$ (blue tit) and $0.654$ (Eucalyptus). Our
conditions range from $0.002$ to $0.551$, and the low end is extreme:
\texttt{bt\_sol\_naive} sits at $0.005$, meaning 38 analyses with effectively
one error bar between them.

Branching raises this but does not close it. \texttt{bt\_sol\_naive} $0.005 \to$
\texttt{bt\_sol\_continuation} $0.088 \to$ \texttt{bt\_sol\_spread} $0.144$,
against a human $0.584$. Human analysts differ in how much information their
model extracts from the data; ours mostly do not.

\subsection{Where the arms sit relative to the crowd is dataset-dependent}
\label{sec:results-location}

On blue tit the arms never reach the weak-effect region the crowd occupies. 37
of 111 human analyses (33\%) report $|t| < 4$; across all nine blue-tit
conditions, 8 of 282 in-scope analyses do, and exactly one crosses zero against
the crowd's seven. The model talks itself into a confident negative effect
nearly every time.

On Eucalyptus the same claim is false, and we state that rather than generalize
from the first dataset. There the true effect is near zero, 80 of 85 human
analyses (94\%) report $|t| < 4$, and so do 434 of our 438; 62 of those 438
cross zero, including 21 of 49 in \texttt{eu\_sol\_continuation}. Where
the data support a null, the arms find it.

The honest summary is that homogeneity of \emph{precision}
(Section~\ref{sec:results-homogeneity}) replicates across both datasets, while
failure to reach the null region does not: it is a blue-tit finding, and on that
dataset it is the largest gap to the human panel that no sampler or
standardization change addresses.

\subsection{Attrition, and what it is selecting on}
\label{sec:results-attrition}

Eucalyptus attrition is now negligible --- 0 to 2 analyses lost per condition of
50 --- after the degrees-of-freedom policy in the setup was made per-dataset.
Before it, the same runs lost 50--90\%, and the survivors were selected by
whether the analyst wrote ``$t$ with df'' or ``asymptotic $z$'': a reporting
convention, not an analytic property.

Blue tit remains uneven. \texttt{bt\_opus\_naive} retains 10 of 50 and
\texttt{bt\_opus\_continuation} 19 of 49, against 8--12 lost for
\texttt{gpt-5.6-sol}. Those rows report a cluster-robust $z$ with no stated
cluster count, and on a dataset whose comparator standardizes on the cluster
scale we decline to guess one. The consequence is that the two blue-tit
\texttt{claude-opus-5} columns rest on small samples and should not carry weight
in a model comparison.

\subsection{A diagnostic on the raw test statistic}
\label{sec:results-teststat}

Because Layer 2 is the only place a common scale is imposed, it is worth asking
whether it creates or destroys the arm differences. The raw focal statistic
answers that: it consumes no degrees of freedom and no standardization rule.

Its validity as a proxy depends on the dataset. Eucalyptus analyses are fit on
an essentially fixed sample ($n_{\text{obs}}$ coefficient of variation $0.063$),
and there the rank correlation between $|t|$ and the standardized $|Z_r|$ is
$\rho = +0.999$: an arm contrast on the raw statistic is a genuine robustness
check. Blue-tit analyses are not, because they sit at both chick and nest level
($n_{\text{obs}}$ from 453 to 3{,}635, coefficient of variation $0.528$), and
$\rho$ falls to $+0.621$.

Neither should be compared to the humans on this scale. Crowd $|t|$ spans
0--105 (blue tit) and 0--467 (Eucalyptus) in part because crowd sample sizes
vary far more than ours. The test statistic is an internal diagnostic; Fisher
$z$ remains the only scale on which the arms and the human panel are comparable.

\subsection{Threats to validity}
\label{sec:results-threats}

\begin{itemize}
\item \textbf{Widths are lower bounds.} No condition converged
(Section~\ref{sec:results-convergence}), so every spread figure is a partial
sample of the decision space.
\item \textbf{Spread is measured on a 95\% window.} At $n \approx 50$ that is
close to the raw range and is set by the two most extreme analyses. Blue tit /
\texttt{gpt-5.6-sol} Branchpoint reads 56\% of the human panel on this window
against 37\% on a 90\% one. Coverage does not have that sensitivity.
\item \textbf{One seed per condition.} The grid is seed 7 throughout. A
three-seed replication of one cell gave widths of $0.19$, $0.40$ and $0.45$ ---
seed variance is of the same order as several of the arm differences reported
here, and single-cell comparisons should not be read as significant.
\item \textbf{Ban lists are per-trajectory.} The same decision axis reached from
a different parent starts with an empty ban list, so distinct branches can and
do land on identical specifications. This deflates coverage by an amount we have
not quantified.
\item \textbf{BpDynamic ran for one model.} Its coverage advantage is consistent
across the two datasets it was run on, but rests on two conditions.
\end{itemize}

\end{document}
; also compiles
% standalone (the guard falls back to article when the workshop style is absent).
\documentclass{article}

\IfFileExists{neurips_2026.sty}{\usepackage[dblblindworkshop]{neurips_2026}}{%
  % standalone fallback: lmodern gives microtype a scalable font to expand
  \usepackage[margin=1in]{geometry}\usepackage{lmodern}}

\usepackage[utf8]{inputenc}
\usepackage[T1]{fontenc}
\usepackage{hyperref}
\usepackage{url}
\usepackage{booktabs}
\usepackage{amsfonts}
\usepackage{amsmath}
\usepackage{amssymb}
\usepackage{microtype}
\usepackage{graphicx}
\usepackage{xcolor}

\title{Branchpoint: Towards Structured Multiverse Analysis}

\author{%
  Anonymous Author(s) \\
  Affiliation \\
  \texttt{email} \\
}

\begin{document}
\maketitle

\section{Results}
\label{sec:results}

Eighteen conditions, $n = 50$ analyses each: three models $\times$ two datasets
$\times$ \{Naive, Branchpoint, BpDynamic\}. 896 analyses in total, at a compute
cost of \$145.

\subsection{Branching produces distinct analyses; the allocator decides how many}
\label{sec:results-coverage}

% Generated by scripts/make_coverage_table.py -- do not hand-edit.
\begin{table}[t]
\centering
\small
\begin{tabular}{llrr}
\toprule
Sampling & Model & Width & \$/unique \\
\midrule
\multicolumn{4}{l}{\textbf{Blue tit}} \\
\quad Human & --- & 0.98 & --- \\
\addlinespace[2pt]
\quad Naive & gpt-5.6-sol & 0.05 & 0.59 \\
 & claude-opus-5 & 0.45 & 3.50 \\
 & glm-5.3 & 0.53 & 0.11 \\
\addlinespace[2pt]
\quad Branchpoint & gpt-5.6-sol & 0.55 & 0.20 \\
 & claude-opus-5 & 0.46 & 0.53 \\
 & glm-5.3 & 0.54 & 0.17 \\
\addlinespace[2pt]
\quad BpDynamic & gpt-5.6-sol & 0.30 & 0.18 \\
 & claude-opus-5 & 0.46 & 0.47 \\
 & glm-5.3 & 0.55 & 0.11 \\
\addlinespace[2pt]
\midrule
\multicolumn{4}{l}{\textbf{Eucalyptus}} \\
\quad Human & --- & 0.79 & --- \\
\addlinespace[2pt]
\quad Naive & gpt-5.6-sol & 0.03 & 1.18 \\
 & claude-opus-5 & 0.02 & 1.28 \\
 & glm-5.3 & 0.11 & 0.16 \\
\addlinespace[2pt]
\quad Branchpoint & gpt-5.6-sol & 0.10 & 0.25 \\
 & claude-opus-5 & 0.10 & 0.53 \\
 & glm-5.3 & 0.10 & 0.09 \\
\addlinespace[2pt]
\quad BpDynamic & gpt-5.6-sol & 0.12 & 0.21 \\
 & claude-opus-5 & 0.30 & 0.38 \\
 & glm-5.3 & 0.24 & 0.14 \\
\addlinespace[2pt]
\bottomrule
\end{tabular}
\caption{Width and cost by sampling arm, $n=50$ per condition. \emph{Width} is the 95\% window of the Fisher-$z$ distribution, in absolute units, so the human panel is directly comparable. \emph{\$/unique} is total spend for the condition, Layer 1 and its Layer-2 judging together, divided by the number of distinct endpoints it bought; the human panel has no comparable figure. Against Branchpoint on the same cell, BpDynamic is wider in 4 of 6; on Eucalyptus the widths separate the arms, on blue tit they do not.}
\label{tab:coverage-spread}
\end{table}


Table~\ref{tab:coverage-spread} reports the 95\% width of each condition's
Fisher-$z$ distribution in absolute units, against the human panel for the same
dataset, together with what each distinct answer cost.

A second quantity, not tabulated, is needed to read those widths.
\emph{Coverage} is the share of in-scope analyses that are a distinct
standardized endpoint --- how much of what an arm produced is a different answer
rather than the same one again. It rises across the arms for every model and
both datasets: on \texttt{gpt-5.6-sol} it goes $39 \to 68 \to 95$ on blue tit
and $20 \to 69 \to 80$ on Eucalyptus. The naive controls show why it matters.
\texttt{eu\_opus\_naive} produced 50 analyses containing 10 distinct answers,
and \texttt{bt\_opus\_naive} 10 in-scope analyses containing 3. Width alone
reports those as narrow; coverage reports them as \emph{repetitive}, which is a
different failure and the one branching is meant to address.

The allocator does what it is designed to do, and the effect is exact rather
than statistical. The BpDynamic conditions spent 50 branches on 50 distinct
decisions. The Branchpoint conditions, drawing at random from the same open set,
spent 50 branches on 34--37 distinct decisions, wasting the remainder on
decisions already branched.

\paragraph{Covering more decisions does not widen the distribution.}
This is the result we would least like to be true. Across the six cells where
both arms ran, BpDynamic covers more of the decision space in four and is wider
in four --- but they are not the same four. Only \texttt{eu}/\texttt{sol} and
\texttt{eu}/\texttt{opus} win on both. On blue tit BpDynamic leads on coverage
in two of three cells and is the narrower arm in two of three. The additional
decisions the allocator reaches are, on that dataset, ones whose alternatives do
not move the estimand.

\subsection{Nothing converged}
\label{sec:results-convergence}

Every branching condition stopped at \texttt{cap-reached}. Across the twelve of
them, 2{,}157 of 2{,}159 detected decisions remain open: \emph{two} were ever
explored to exhaustion. With $n = 50$ against roughly 200 detected decisions per
condition, even perfect allocation reaches a quarter of them at one alternative
apiece.

Every width reported here is therefore a lower bound from a partial sample, and
the convergence curves are still climbing at $n = 50$. This is why coverage is
worth stating alongside the widths: a width without a denominator invites
reading a 25\%-coverage sample as an estimate of the multiverse.

\subsection{The analyses are far more homogeneous than the humans'}
\label{sec:results-homogeneity}

Under the Fisher-$z$ convention the standard error depends only on the degrees of
freedom, so its variation across a condition measures how much the analyses
differ in effective sample size. The human panels have a standard-error
coefficient of variation of $0.584$ (blue tit) and $0.654$ (Eucalyptus). Our
conditions range from $0.002$ to $0.551$, and the low end is extreme:
\texttt{bt\_sol\_naive} sits at $0.005$, meaning 38 analyses with effectively
one error bar between them.

Branching raises this but does not close it. \texttt{bt\_sol\_naive} $0.005 \to$
\texttt{bt\_sol\_continuation} $0.088 \to$ \texttt{bt\_sol\_spread} $0.144$,
against a human $0.584$. Human analysts differ in how much information their
model extracts from the data; ours mostly do not.

\subsection{Where the arms sit relative to the crowd is dataset-dependent}
\label{sec:results-location}

On blue tit the arms never reach the weak-effect region the crowd occupies. 37
of 111 human analyses (33\%) report $|t| < 4$; across all nine blue-tit
conditions, 8 of 282 in-scope analyses do, and exactly one crosses zero against
the crowd's seven. The model talks itself into a confident negative effect
nearly every time.

On Eucalyptus the same claim is false, and we state that rather than generalize
from the first dataset. There the true effect is near zero, 80 of 85 human
analyses (94\%) report $|t| < 4$, and so do 434 of our 438; 62 of those 438
cross zero, including 21 of 49 in \texttt{eu\_sol\_continuation}. Where
the data support a null, the arms find it.

The honest summary is that homogeneity of \emph{precision}
(Section~\ref{sec:results-homogeneity}) replicates across both datasets, while
failure to reach the null region does not: it is a blue-tit finding, and on that
dataset it is the largest gap to the human panel that no sampler or
standardization change addresses.

\subsection{Attrition, and what it is selecting on}
\label{sec:results-attrition}

Eucalyptus attrition is now negligible --- 0 to 2 analyses lost per condition of
50 --- after the degrees-of-freedom policy in the setup was made per-dataset.
Before it, the same runs lost 50--90\%, and the survivors were selected by
whether the analyst wrote ``$t$ with df'' or ``asymptotic $z$'': a reporting
convention, not an analytic property.

Blue tit remains uneven. \texttt{bt\_opus\_naive} retains 10 of 50 and
\texttt{bt\_opus\_continuation} 19 of 49, against 8--12 lost for
\texttt{gpt-5.6-sol}. Those rows report a cluster-robust $z$ with no stated
cluster count, and on a dataset whose comparator standardizes on the cluster
scale we decline to guess one. The consequence is that the two blue-tit
\texttt{claude-opus-5} columns rest on small samples and should not carry weight
in a model comparison.

\subsection{A diagnostic on the raw test statistic}
\label{sec:results-teststat}

Because Layer 2 is the only place a common scale is imposed, it is worth asking
whether it creates or destroys the arm differences. The raw focal statistic
answers that: it consumes no degrees of freedom and no standardization rule.

Its validity as a proxy depends on the dataset. Eucalyptus analyses are fit on
an essentially fixed sample ($n_{\text{obs}}$ coefficient of variation $0.063$),
and there the rank correlation between $|t|$ and the standardized $|Z_r|$ is
$\rho = +0.999$: an arm contrast on the raw statistic is a genuine robustness
check. Blue-tit analyses are not, because they sit at both chick and nest level
($n_{\text{obs}}$ from 453 to 3{,}635, coefficient of variation $0.528$), and
$\rho$ falls to $+0.621$.

Neither should be compared to the humans on this scale. Crowd $|t|$ spans
0--105 (blue tit) and 0--467 (Eucalyptus) in part because crowd sample sizes
vary far more than ours. The test statistic is an internal diagnostic; Fisher
$z$ remains the only scale on which the arms and the human panel are comparable.

\subsection{Threats to validity}
\label{sec:results-threats}

\begin{itemize}
\item \textbf{Widths are lower bounds.} No condition converged
(Section~\ref{sec:results-convergence}), so every spread figure is a partial
sample of the decision space.
\item \textbf{Spread is measured on a 95\% window.} At $n \approx 50$ that is
close to the raw range and is set by the two most extreme analyses. Blue tit /
\texttt{gpt-5.6-sol} Branchpoint reads 56\% of the human panel on this window
against 37\% on a 90\% one. Coverage does not have that sensitivity.
\item \textbf{One seed per condition.} The grid is seed 7 throughout. A
three-seed replication of one cell gave widths of $0.19$, $0.40$ and $0.45$ ---
seed variance is of the same order as several of the arm differences reported
here, and single-cell comparisons should not be read as significant.
\item \textbf{Ban lists are per-trajectory.} The same decision axis reached from
a different parent starts with an empty ban list, so distinct branches can and
do land on identical specifications. This deflates coverage by an amount we have
not quantified.
\item \textbf{BpDynamic ran for one model.} Its coverage advantage is consistent
across the two datasets it was run on, but rests on two conditions.
\end{itemize}

\end{document}
; also compiles
% standalone (the guard falls back to article when the workshop style is absent).
\documentclass{article}

\IfFileExists{neurips_2026.sty}{\usepackage[dblblindworkshop]{neurips_2026}}{%
  % standalone fallback: lmodern gives microtype a scalable font to expand
  \usepackage[margin=1in]{geometry}\usepackage{lmodern}}

\usepackage[utf8]{inputenc}
\usepackage[T1]{fontenc}
\usepackage{hyperref}
\usepackage{url}
\usepackage{booktabs}
\usepackage{amsfonts}
\usepackage{amsmath}
\usepackage{amssymb}
\usepackage{microtype}
\usepackage{graphicx}
\usepackage{xcolor}

\title{Branchpoint: Towards Structured Multiverse Analysis}

\author{%
  Anonymous Author(s) \\
  Affiliation \\
  \texttt{email} \\
}

\begin{document}
\maketitle

\section{Results}
\label{sec:results}

Eighteen conditions, $n = 50$ analyses each: three models $\times$ two datasets
$\times$ \{Naive, Branchpoint, BpDynamic\}. 896 analyses in total, at a compute
cost of \$145.

\subsection{Branching produces distinct analyses; the allocator decides how many}
\label{sec:results-coverage}

% Generated by scripts/make_coverage_table.py -- do not hand-edit.
\begin{table}[t]
\centering
\small
\begin{tabular}{llrr}
\toprule
Sampling & Model & Width & \$/unique \\
\midrule
\multicolumn{4}{l}{\textbf{Blue tit}} \\
\quad Human & --- & 0.98 & --- \\
\addlinespace[2pt]
\quad Naive & gpt-5.6-sol & 0.05 & 0.59 \\
 & claude-opus-5 & 0.45 & 3.50 \\
 & glm-5.3 & 0.53 & 0.11 \\
\addlinespace[2pt]
\quad Branchpoint & gpt-5.6-sol & 0.55 & 0.20 \\
 & claude-opus-5 & 0.46 & 0.53 \\
 & glm-5.3 & 0.54 & 0.17 \\
\addlinespace[2pt]
\quad BpDynamic & gpt-5.6-sol & 0.30 & 0.18 \\
 & claude-opus-5 & 0.46 & 0.47 \\
 & glm-5.3 & 0.55 & 0.11 \\
\addlinespace[2pt]
\midrule
\multicolumn{4}{l}{\textbf{Eucalyptus}} \\
\quad Human & --- & 0.79 & --- \\
\addlinespace[2pt]
\quad Naive & gpt-5.6-sol & 0.03 & 1.18 \\
 & claude-opus-5 & 0.02 & 1.28 \\
 & glm-5.3 & 0.11 & 0.16 \\
\addlinespace[2pt]
\quad Branchpoint & gpt-5.6-sol & 0.10 & 0.25 \\
 & claude-opus-5 & 0.10 & 0.53 \\
 & glm-5.3 & 0.10 & 0.09 \\
\addlinespace[2pt]
\quad BpDynamic & gpt-5.6-sol & 0.12 & 0.21 \\
 & claude-opus-5 & 0.30 & 0.38 \\
 & glm-5.3 & 0.24 & 0.14 \\
\addlinespace[2pt]
\bottomrule
\end{tabular}
\caption{Width and cost by sampling arm, $n=50$ per condition. \emph{Width} is the 95\% window of the Fisher-$z$ distribution, in absolute units, so the human panel is directly comparable. \emph{\$/unique} is total spend for the condition, Layer 1 and its Layer-2 judging together, divided by the number of distinct endpoints it bought; the human panel has no comparable figure. Against Branchpoint on the same cell, BpDynamic is wider in 4 of 6; on Eucalyptus the widths separate the arms, on blue tit they do not.}
\label{tab:coverage-spread}
\end{table}


Table~\ref{tab:coverage-spread} reports the 95\% width of each condition's
Fisher-$z$ distribution in absolute units, against the human panel for the same
dataset, together with what each distinct answer cost.

A second quantity, not tabulated, is needed to read those widths.
\emph{Coverage} is the share of in-scope analyses that are a distinct
standardized endpoint --- how much of what an arm produced is a different answer
rather than the same one again. It rises across the arms for every model and
both datasets: on \texttt{gpt-5.6-sol} it goes $39 \to 68 \to 95$ on blue tit
and $20 \to 69 \to 80$ on Eucalyptus. The naive controls show why it matters.
\texttt{eu\_opus\_naive} produced 50 analyses containing 10 distinct answers,
and \texttt{bt\_opus\_naive} 10 in-scope analyses containing 3. Width alone
reports those as narrow; coverage reports them as \emph{repetitive}, which is a
different failure and the one branching is meant to address.

The allocator does what it is designed to do, and the effect is exact rather
than statistical. The BpDynamic conditions spent 50 branches on 50 distinct
decisions. The Branchpoint conditions, drawing at random from the same open set,
spent 50 branches on 34--37 distinct decisions, wasting the remainder on
decisions already branched.

\paragraph{Covering more decisions does not widen the distribution.}
This is the result we would least like to be true. Across the six cells where
both arms ran, BpDynamic covers more of the decision space in four and is wider
in four --- but they are not the same four. Only \texttt{eu}/\texttt{sol} and
\texttt{eu}/\texttt{opus} win on both. On blue tit BpDynamic leads on coverage
in two of three cells and is the narrower arm in two of three. The additional
decisions the allocator reaches are, on that dataset, ones whose alternatives do
not move the estimand.

\subsection{Nothing converged}
\label{sec:results-convergence}

Every branching condition stopped at \texttt{cap-reached}. Across the twelve of
them, 2{,}157 of 2{,}159 detected decisions remain open: \emph{two} were ever
explored to exhaustion. With $n = 50$ against roughly 200 detected decisions per
condition, even perfect allocation reaches a quarter of them at one alternative
apiece.

Every width reported here is therefore a lower bound from a partial sample, and
the convergence curves are still climbing at $n = 50$. This is why coverage is
worth stating alongside the widths: a width without a denominator invites
reading a 25\%-coverage sample as an estimate of the multiverse.

\subsection{The analyses are far more homogeneous than the humans'}
\label{sec:results-homogeneity}

Under the Fisher-$z$ convention the standard error depends only on the degrees of
freedom, so its variation across a condition measures how much the analyses
differ in effective sample size. The human panels have a standard-error
coefficient of variation of $0.584$ (blue tit) and $0.654$ (Eucalyptus). Our
conditions range from $0.002$ to $0.551$, and the low end is extreme:
\texttt{bt\_sol\_naive} sits at $0.005$, meaning 38 analyses with effectively
one error bar between them.

Branching raises this but does not close it. \texttt{bt\_sol\_naive} $0.005 \to$
\texttt{bt\_sol\_continuation} $0.088 \to$ \texttt{bt\_sol\_spread} $0.144$,
against a human $0.584$. Human analysts differ in how much information their
model extracts from the data; ours mostly do not.

\subsection{Where the arms sit relative to the crowd is dataset-dependent}
\label{sec:results-location}

On blue tit the arms never reach the weak-effect region the crowd occupies. 37
of 111 human analyses (33\%) report $|t| < 4$; across all nine blue-tit
conditions, 8 of 282 in-scope analyses do, and exactly one crosses zero against
the crowd's seven. The model talks itself into a confident negative effect
nearly every time.

On Eucalyptus the same claim is false, and we state that rather than generalize
from the first dataset. There the true effect is near zero, 80 of 85 human
analyses (94\%) report $|t| < 4$, and so do 434 of our 438; 62 of those 438
cross zero, including 21 of 49 in \texttt{eu\_sol\_continuation}. Where
the data support a null, the arms find it.

The honest summary is that homogeneity of \emph{precision}
(Section~\ref{sec:results-homogeneity}) replicates across both datasets, while
failure to reach the null region does not: it is a blue-tit finding, and on that
dataset it is the largest gap to the human panel that no sampler or
standardization change addresses.

\subsection{Attrition, and what it is selecting on}
\label{sec:results-attrition}

Eucalyptus attrition is now negligible --- 0 to 2 analyses lost per condition of
50 --- after the degrees-of-freedom policy in the setup was made per-dataset.
Before it, the same runs lost 50--90\%, and the survivors were selected by
whether the analyst wrote ``$t$ with df'' or ``asymptotic $z$'': a reporting
convention, not an analytic property.

Blue tit remains uneven. \texttt{bt\_opus\_naive} retains 10 of 50 and
\texttt{bt\_opus\_continuation} 19 of 49, against 8--12 lost for
\texttt{gpt-5.6-sol}. Those rows report a cluster-robust $z$ with no stated
cluster count, and on a dataset whose comparator standardizes on the cluster
scale we decline to guess one. The consequence is that the two blue-tit
\texttt{claude-opus-5} columns rest on small samples and should not carry weight
in a model comparison.

\subsection{A diagnostic on the raw test statistic}
\label{sec:results-teststat}

Because Layer 2 is the only place a common scale is imposed, it is worth asking
whether it creates or destroys the arm differences. The raw focal statistic
answers that: it consumes no degrees of freedom and no standardization rule.

Its validity as a proxy depends on the dataset. Eucalyptus analyses are fit on
an essentially fixed sample ($n_{\text{obs}}$ coefficient of variation $0.063$),
and there the rank correlation between $|t|$ and the standardized $|Z_r|$ is
$\rho = +0.999$: an arm contrast on the raw statistic is a genuine robustness
check. Blue-tit analyses are not, because they sit at both chick and nest level
($n_{\text{obs}}$ from 453 to 3{,}635, coefficient of variation $0.528$), and
$\rho$ falls to $+0.621$.

Neither should be compared to the humans on this scale. Crowd $|t|$ spans
0--105 (blue tit) and 0--467 (Eucalyptus) in part because crowd sample sizes
vary far more than ours. The test statistic is an internal diagnostic; Fisher
$z$ remains the only scale on which the arms and the human panel are comparable.

\subsection{Threats to validity}
\label{sec:results-threats}

\begin{itemize}
\item \textbf{Widths are lower bounds.} No condition converged
(Section~\ref{sec:results-convergence}), so every spread figure is a partial
sample of the decision space.
\item \textbf{Spread is measured on a 95\% window.} At $n \approx 50$ that is
close to the raw range and is set by the two most extreme analyses. Blue tit /
\texttt{gpt-5.6-sol} Branchpoint reads 56\% of the human panel on this window
against 37\% on a 90\% one. Coverage does not have that sensitivity.
\item \textbf{One seed per condition.} The grid is seed 7 throughout. A
three-seed replication of one cell gave widths of $0.19$, $0.40$ and $0.45$ ---
seed variance is of the same order as several of the arm differences reported
here, and single-cell comparisons should not be read as significant.
\item \textbf{Ban lists are per-trajectory.} The same decision axis reached from
a different parent starts with an empty ban list, so distinct branches can and
do land on identical specifications. This deflates coverage by an amount we have
not quantified.
\item \textbf{BpDynamic ran for one model.} Its coverage advantage is consistent
across the two datasets it was run on, but rests on two conditions.
\end{itemize}

\end{document}
; also compiles
% standalone (the guard falls back to article when the workshop style is absent).
\documentclass{article}

\IfFileExists{neurips_2026.sty}{\usepackage[dblblindworkshop]{neurips_2026}}{%
  % standalone fallback: lmodern gives microtype a scalable font to expand
  \usepackage[margin=1in]{geometry}\usepackage{lmodern}}

\usepackage[utf8]{inputenc}
\usepackage[T1]{fontenc}
\usepackage{hyperref}
\usepackage{url}
\usepackage{booktabs}
\usepackage{amsfonts}
\usepackage{amsmath}
\usepackage{amssymb}
\usepackage{microtype}
\usepackage{graphicx}
\usepackage{xcolor}

\title{Branchpoint: Towards Structured Multiverse Analysis}

\author{%
  Anonymous Author(s) \\
  Affiliation \\
  \texttt{email} \\
}

\begin{document}
\maketitle

\section{Results}
\label{sec:results}

Eighteen conditions, $n = 50$ analyses each: three models $\times$ two datasets
$\times$ \{Naive, Branchpoint, BpDynamic\}. 896 analyses in total, at a compute
cost of \$145.

\subsection{Branching produces distinct analyses; the allocator decides how many}
\label{sec:results-coverage}

% Generated by scripts/make_coverage_table.py -- do not hand-edit.
\begin{table}[t]
\centering
\small
\begin{tabular}{llrr}
\toprule
Sampling & Model & Width & \$/unique \\
\midrule
\multicolumn{4}{l}{\textbf{Blue tit}} \\
\quad Human & --- & 0.98 & --- \\
\addlinespace[2pt]
\quad Naive & gpt-5.6-sol & 0.05 & 0.59 \\
 & claude-opus-5 & 0.45 & 3.50 \\
 & glm-5.3 & 0.53 & 0.11 \\
\addlinespace[2pt]
\quad Branchpoint & gpt-5.6-sol & 0.55 & 0.20 \\
 & claude-opus-5 & 0.46 & 0.53 \\
 & glm-5.3 & 0.54 & 0.17 \\
\addlinespace[2pt]
\quad BpDynamic & gpt-5.6-sol & 0.30 & 0.18 \\
 & claude-opus-5 & 0.46 & 0.47 \\
 & glm-5.3 & 0.55 & 0.11 \\
\addlinespace[2pt]
\midrule
\multicolumn{4}{l}{\textbf{Eucalyptus}} \\
\quad Human & --- & 0.79 & --- \\
\addlinespace[2pt]
\quad Naive & gpt-5.6-sol & 0.03 & 1.18 \\
 & claude-opus-5 & 0.02 & 1.28 \\
 & glm-5.3 & 0.11 & 0.16 \\
\addlinespace[2pt]
\quad Branchpoint & gpt-5.6-sol & 0.10 & 0.25 \\
 & claude-opus-5 & 0.10 & 0.53 \\
 & glm-5.3 & 0.10 & 0.09 \\
\addlinespace[2pt]
\quad BpDynamic & gpt-5.6-sol & 0.12 & 0.21 \\
 & claude-opus-5 & 0.30 & 0.38 \\
 & glm-5.3 & 0.24 & 0.14 \\
\addlinespace[2pt]
\bottomrule
\end{tabular}
\caption{Width and cost by sampling arm, $n=50$ per condition. \emph{Width} is the 95\% window of the Fisher-$z$ distribution, in absolute units, so the human panel is directly comparable. \emph{\$/unique} is total spend for the condition, Layer 1 and its Layer-2 judging together, divided by the number of distinct endpoints it bought; the human panel has no comparable figure. Against Branchpoint on the same cell, BpDynamic is wider in 4 of 6; on Eucalyptus the widths separate the arms, on blue tit they do not.}
\label{tab:coverage-spread}
\end{table}


Table~\ref{tab:coverage-spread} reports the 95\% width of each condition's
Fisher-$z$ distribution in absolute units, against the human panel for the same
dataset, together with what each distinct answer cost.

A second quantity, not tabulated, is needed to read those widths.
\emph{Coverage} is the share of in-scope analyses that are a distinct
standardized endpoint --- how much of what an arm produced is a different answer
rather than the same one again. It rises across the arms for every model and
both datasets: on \texttt{gpt-5.6-sol} it goes $39 \to 68 \to 95$ on blue tit
and $20 \to 69 \to 80$ on Eucalyptus. The naive controls show why it matters.
\texttt{eu\_opus\_naive} produced 50 analyses containing 10 distinct answers,
and \texttt{bt\_opus\_naive} 10 in-scope analyses containing 3. Width alone
reports those as narrow; coverage reports them as \emph{repetitive}, which is a
different failure and the one branching is meant to address.

The allocator does what it is designed to do, and the effect is exact rather
than statistical. The BpDynamic conditions spent 50 branches on 50 distinct
decisions. The Branchpoint conditions, drawing at random from the same open set,
spent 50 branches on 34--37 distinct decisions, wasting the remainder on
decisions already branched.

\paragraph{Covering more decisions does not widen the distribution.}
This is the result we would least like to be true. Across the six cells where
both arms ran, BpDynamic covers more of the decision space in four and is wider
in four --- but they are not the same four. Only \texttt{eu}/\texttt{sol} and
\texttt{eu}/\texttt{opus} win on both. On blue tit BpDynamic leads on coverage
in two of three cells and is the narrower arm in two of three. The additional
decisions the allocator reaches are, on that dataset, ones whose alternatives do
not move the estimand.

\subsection{Nothing converged}
\label{sec:results-convergence}

Every branching condition stopped at \texttt{cap-reached}. Across the twelve of
them, 2{,}157 of 2{,}159 detected decisions remain open: \emph{two} were ever
explored to exhaustion. With $n = 50$ against roughly 200 detected decisions per
condition, even perfect allocation reaches a quarter of them at one alternative
apiece.

Every width reported here is therefore a lower bound from a partial sample, and
the convergence curves are still climbing at $n = 50$. This is why coverage is
worth stating alongside the widths: a width without a denominator invites
reading a 25\%-coverage sample as an estimate of the multiverse.

\subsection{The analyses are far more homogeneous than the humans'}
\label{sec:results-homogeneity}

Under the Fisher-$z$ convention the standard error depends only on the degrees of
freedom, so its variation across a condition measures how much the analyses
differ in effective sample size. The human panels have a standard-error
coefficient of variation of $0.584$ (blue tit) and $0.654$ (Eucalyptus). Our
conditions range from $0.002$ to $0.551$, and the low end is extreme:
\texttt{bt\_sol\_naive} sits at $0.005$, meaning 38 analyses with effectively
one error bar between them.

Branching raises this but does not close it. \texttt{bt\_sol\_naive} $0.005 \to$
\texttt{bt\_sol\_continuation} $0.088 \to$ \texttt{bt\_sol\_spread} $0.144$,
against a human $0.584$. Human analysts differ in how much information their
model extracts from the data; ours mostly do not.

\subsection{Where the arms sit relative to the crowd is dataset-dependent}
\label{sec:results-location}

On blue tit the arms never reach the weak-effect region the crowd occupies. 37
of 111 human analyses (33\%) report $|t| < 4$; across all nine blue-tit
conditions, 8 of 282 in-scope analyses do, and exactly one crosses zero against
the crowd's seven. The model talks itself into a confident negative effect
nearly every time.

On Eucalyptus the same claim is false, and we state that rather than generalize
from the first dataset. There the true effect is near zero, 80 of 85 human
analyses (94\%) report $|t| < 4$, and so do 434 of our 438; 62 of those 438
cross zero, including 21 of 49 in \texttt{eu\_sol\_continuation}. Where
the data support a null, the arms find it.

The honest summary is that homogeneity of \emph{precision}
(Section~\ref{sec:results-homogeneity}) replicates across both datasets, while
failure to reach the null region does not: it is a blue-tit finding, and on that
dataset it is the largest gap to the human panel that no sampler or
standardization change addresses.

\subsection{Attrition, and what it is selecting on}
\label{sec:results-attrition}

Eucalyptus attrition is now negligible --- 0 to 2 analyses lost per condition of
50 --- after the degrees-of-freedom policy in the setup was made per-dataset.
Before it, the same runs lost 50--90\%, and the survivors were selected by
whether the analyst wrote ``$t$ with df'' or ``asymptotic $z$'': a reporting
convention, not an analytic property.

Blue tit remains uneven. \texttt{bt\_opus\_naive} retains 10 of 50 and
\texttt{bt\_opus\_continuation} 19 of 49, against 8--12 lost for
\texttt{gpt-5.6-sol}. Those rows report a cluster-robust $z$ with no stated
cluster count, and on a dataset whose comparator standardizes on the cluster
scale we decline to guess one. The consequence is that the two blue-tit
\texttt{claude-opus-5} columns rest on small samples and should not carry weight
in a model comparison.

\subsection{A diagnostic on the raw test statistic}
\label{sec:results-teststat}

Because Layer 2 is the only place a common scale is imposed, it is worth asking
whether it creates or destroys the arm differences. The raw focal statistic
answers that: it consumes no degrees of freedom and no standardization rule.

Its validity as a proxy depends on the dataset. Eucalyptus analyses are fit on
an essentially fixed sample ($n_{\text{obs}}$ coefficient of variation $0.063$),
and there the rank correlation between $|t|$ and the standardized $|Z_r|$ is
$\rho = +0.999$: an arm contrast on the raw statistic is a genuine robustness
check. Blue-tit analyses are not, because they sit at both chick and nest level
($n_{\text{obs}}$ from 453 to 3{,}635, coefficient of variation $0.528$), and
$\rho$ falls to $+0.621$.

Neither should be compared to the humans on this scale. Crowd $|t|$ spans
0--105 (blue tit) and 0--467 (Eucalyptus) in part because crowd sample sizes
vary far more than ours. The test statistic is an internal diagnostic; Fisher
$z$ remains the only scale on which the arms and the human panel are comparable.

\subsection{Threats to validity}
\label{sec:results-threats}

\begin{itemize}
\item \textbf{Widths are lower bounds.} No condition converged
(Section~\ref{sec:results-convergence}), so every spread figure is a partial
sample of the decision space.
\item \textbf{Spread is measured on a 95\% window.} At $n \approx 50$ that is
close to the raw range and is set by the two most extreme analyses. Blue tit /
\texttt{gpt-5.6-sol} Branchpoint reads 56\% of the human panel on this window
against 37\% on a 90\% one. Coverage does not have that sensitivity.
\item \textbf{One seed per condition.} The grid is seed 7 throughout. A
three-seed replication of one cell gave widths of $0.19$, $0.40$ and $0.45$ ---
seed variance is of the same order as several of the arm differences reported
here, and single-cell comparisons should not be read as significant.
\item \textbf{Ban lists are per-trajectory.} The same decision axis reached from
a different parent starts with an empty ban list, so distinct branches can and
do land on identical specifications. This deflates coverage by an amount we have
not quantified.
\item \textbf{BpDynamic ran for one model.} Its coverage advantage is consistent
across the two datasets it was run on, but rests on two conditions.
\end{itemize}

\end{document}
